\documentclass[11pt]{article}
\usepackage[T1]{fontenc}
\usepackage[utf8]{inputenc}
\usepackage[margin=1in]{geometry}
\usepackage{amsmath,amssymb,amsthm}
\usepackage{microtype}
\usepackage[hidelinks]{hyperref}
\newcommand{\E}{\mathbb E}
\newcommand{\Pb}{\mathbb P}
\newcommand{\Ind}[1]{\mathbf 1\{#1\}}
\newcommand{\Bfun}{\mathcal B}
\DeclareMathOperator*{\argmax}{arg\,max}
\allowdisplaybreaks[1]
\title{Two Strategic Sellers without Forgetting}
\author{Analytical progress note}
\date{October 3, 2026}

\begin{document}
\maketitle

This note extends the paper's cumulative-learning model by replacing the
known outside option with a second long-lived strategic seller. It derives
customer demand, state transitions, policy evaluation, best responses, and
the equilibrium product-choice conditions. A finite-horizon benchmark is
specified analytically; no numerical experiments are undertaken.

\paragraph{Notation.}
We retain the paper's $p_i,c_i,R,\gamma$, its selection, product, and outcome
variables $y_t,i_t,u_t$, and its discounted continuation gap $G$.
Following the previous progress notes, $U$ denotes a value under a fixed
policy and $V$ an optimized value. Thus $U_j^\sigma$ will evaluate a fixed
policy pair, whereas $V_j^{\sigma_{-j}}$ will denote a best-response value.
The pointwise product-choice gain is $a_j$, extending $a_\rho$ in the
previous notes. The symbol $\pi_t^j$ remains period profit; a deterministic
product rule will be denoted by $i_j^H$ in the finite-horizon benchmark.
Seller subscripts distinguish the two copies of the original objects.

\section{Environment and information}

There is one customer and two long-lived strategic sellers, $A$ and $B$.
They have the same products and discount factor:
\begin{equation}\label{eq:primitives}
0<p_1<p_2\leq1,\qquad 0\leq c_1<c_2<R,\qquad \gamma\in(0,1).
\end{equation}
Product $i\in\{1,2\}$ costs $c_i$ and generates success with actual
probability $p_i$. The selected seller receives revenue $R$ and incurs
the cost of its chosen product, regardless of the realized outcome.
The unselected seller receives zero and incurs no product cost.

As in the paper, $y_t\in\{A,B\}$ is the seller selected in period $t$,
$i_t\in\{1,2\}$ is that seller's hidden product choice, and
$u_t\in\{0,1\}$ is the outcome, with $u_t=1$ representing success.
Periods are indexed by $t=1,2,\ldots$. Seller $j$'s period profit is
\begin{equation}\label{eq:profit}
\pi_t^j=\Ind{y_t=j}(R-c_{i_t}),\qquad j\in\{A,B\},
\end{equation}
where $\Ind{\cdot}$ is the indicator of an event. There is only one
realized product choice $i_t$ in each period.

Both sellers observe the customer's selections and the outcomes generated
by either seller. They know the primitives and correctly anticipate each
other's equilibrium policy functions. Product choice is private, but the
joint observation state is public to both sellers.

The customer observes outcomes and their source, but does not observe the
selected seller's product choice. She follows a specified learning and
choice rule: separately for each seller $j$, she treats outcomes as if
they were independent Bernoulli observations with a fixed unknown success
probability $\theta_j$. This is her statistical representation of that
seller's perceived performance or reputation. It is not a claim that the
seller's actual success probability is constant under its strategy.
She does not infer the sellers' strategies or reconstruct hidden product
choices from equilibrium behavior.

\section{Timing, state, and customer beliefs}

\subsection{Timing within a period}

The order of events is as follows.
\begin{enumerate}
\item The customer independently draws one sample from her current belief
about each seller.
\item She selects the seller with the higher sample.
\item Only the selected seller privately chooses product $1$ or $2$.
\item An outcome is generated with the chosen product's actual success
probability:
\begin{equation}\label{eq:actual-outcome}
\Pb(u_t=1\mid y_t=j,i_t=i,\text{past observations})=p_i.
\end{equation}
\item The selected seller's success or failure count increases by one.
\end{enumerate}
The customer's samples determine selection only. They are not the
probabilities used to generate $u_t$. In particular, sellers do not pay
for a product before the customer chooses a seller.

\subsection{Cumulative counts and perceived parameters}

For seller $j$, let $S_{j,t}$ and $F_{j,t}$ count the successes and
failures observed before period $t$:
\begin{equation}\label{eq:counts}
S_{j,t}=\sum_{\tau=1}^{t-1}\Ind{y_\tau=j,u_\tau=1},
\qquad
F_{j,t}=\sum_{\tau=1}^{t-1}\Ind{y_\tau=j,u_\tau=0}.
\end{equation}
At a generic decision state, suppress the time index and write
\begin{equation}\label{eq:state}
z=(S_A,F_A,S_B,F_B)\in\mathbb N_0^4,
\qquad z_0=(0,0,0,0),
\end{equation}
where $\mathbb N_0$ is the set of nonnegative integers and $z_0$ denotes
the zero-count initial condition. These are ordinary cumulative counts:
there is no forgetting, normalization, or reduction to posterior means.

Starting from independent $\operatorname{Beta}(1,1)$ priors, the customer
maintains independent belief distributions
\begin{equation}\label{eq:beliefs}
\theta_A\mid z\sim\operatorname{Beta}(1+S_A,1+F_A),\qquad
\theta_B\mid z\sim\operatorname{Beta}(1+S_B,1+F_B).
\end{equation}
Indeed, in her representation the likelihood for seller $j$ is
proportional to $\theta_j^{S_j}(1-\theta_j)^{F_j}$; multiplying by its
uniform prior and normalizing gives \eqref{eq:beliefs}. Each new outcome
updates only the belief associated with its source. The perceived
parameters $\theta_j$ and the actual product probabilities $p_i$ have
different roles and must not be identified.

\section{Thompson-sampling demand}

Let $\widetilde\theta_A$ and $\widetilde\theta_B$ be the independent
samples from \eqref{eq:beliefs}. Extending the paper's demand notation,
\begin{equation}\label{eq:demand-definition}
D_A(z)=\Pb(\widetilde\theta_A>\widetilde\theta_B\mid z),
\qquad D_B(z)=1-D_A(z).
\end{equation}
At every finite state both Beta distributions are continuous, so a tie
has probability zero. Both selection probabilities are strictly between
zero and one because the belief densities are positive on $(0,1)$.

For $x\in(0,1)$, define seller $j$'s belief density and cumulative
distribution function by
\begin{equation}\label{eq:density-cdf}
f_j(x\mid z)=\frac{x^{S_j}(1-x)^{F_j}}{\Bfun(1+S_j,1+F_j)},
\qquad H_j(x\mid z)=\int_0^x f_j(b\mid z)\,db.
\end{equation}
Here $\Bfun(\alpha,\beta)=\int_0^1 x^{\alpha-1}(1-x)^{\beta-1}\,dx$
is the Beta function, for positive parameters $\alpha,\beta$.
The letter $H$ for the CDF keeps $F_j$ reserved for failure counts.

Condition first on $A$'s sample being $a$. Independence makes the
conditional chance that $B$'s sample is smaller equal to $H_B(a\mid z)$.
Integrating over $A$'s sample gives
\begin{equation}\label{eq:demand-integral}
\begin{aligned}
D_A(z)
&=\int_0^1\Pb(\widetilde\theta_B<a\mid z)f_A(a\mid z)\,da\\
&=\int_0^1 f_A(a\mid z)H_B(a\mid z)\,da\\
&=\int_0^1\int_0^a f_A(a\mid z)f_B(b\mid z)\,db\,da.
\end{aligned}
\end{equation}
Alternatively, conditional on $B$'s sample being $b$, $A$ wins with
probability $1-H_A(b\mid z)$. Thus the equivalent expression is
\begin{equation}\label{eq:demand-condition-b}
D_A(z)=\int_0^1 f_B(b\mid z)[1-H_A(b\mid z)]\,db.
\end{equation}

\paragraph{Three elementary checks.}
At the initial state both beliefs are uniform. After one success for $A$,
its density is $2a$; after one failure it is $2(1-a)$. In all three cases
$B$'s CDF is $H_B(a)=a$, so
\begin{equation}\label{eq:demand-checks}
\begin{aligned}
D_A(0,0,0,0)&=\int_0^1 a\,da=\frac12,\\
D_A(1,0,0,0)&=\int_0^1 2a^2\,da=\frac23,\\
D_A(0,1,0,0)&=\int_0^1 2(1-a)a\,da=\frac13.
\end{aligned}
\end{equation}
These are direct integral calculations, independent of any seller policy.

\paragraph{Comparison with the known outside option.}
$B$'s posterior mean is
$\E[\theta_B\mid z]=(1+S_B)/(2+S_B+F_B)$, where $\E$ denotes
expectation. Replacing the paper's known $p_0$ by this mean would give
$1-H_A(\E[\theta_B\mid z]\mid z)$. Actual demand in
\eqref{eq:demand-condition-b} is instead the average of
$1-H_A(b\mid z)$ over $B$'s entire belief distribution. These operations
are generally different.

For example, at $z=(2,0,0,0)$, $A$ has a $\operatorname{Beta}(3,1)$
belief and $B$ a uniform belief. Exactly,
\begin{equation*}
D_A(2,0,0,0)=\int_0^1 3a^3\,da=\frac34,
\qquad
1-H_A(1/2\mid z)=1-(1/2)^3=\frac78.
\end{equation*}
Thus even this simple state rules out substitution of the posterior mean.
If $B$'s belief were instead concentrated at a known $p_0$, its draw would
equal $p_0$ with probability one, and the comparison would reduce to
\begin{equation}\label{eq:outside-demand}
\Pb(\widetilde\theta_A>p_0\mid z)
=\int_{p_0}^1 f_A(a\mid z)\,da=D(S_A,F_A),
\end{equation}
the paper's original demand. This recovers the demand comparison; the
two-seller game has a different continuation structure.

At a given state, demand is determined by the customer's two belief
distributions, not directly by the sellers' current product policies.
Product choice takes place after selection and is hidden. Policies affect
future demand through the outcomes they generate and the resulting counts.

\section{State transitions and calendar time}

The superscripts $s$ and $f$ denote success and failure. The four possible
state updates are
\begin{equation}\label{eq:transitions}
\begin{aligned}
T_A^s z&=(S_A+1,F_A,S_B,F_B),\\
T_A^f z&=(S_A,F_A+1,S_B,F_B),\\
T_B^s z&=(S_A,F_A,S_B+1,F_B),\\
T_B^f z&=(S_A,F_A,S_B,F_B+1).
\end{aligned}
\end{equation}
Define the total number of observations and the individual counts by
\begin{equation}\label{eq:layers}
n(z)=S_A+F_A+S_B+F_B,\qquad n_j(z)=S_j+F_j.
\end{equation}
For every $j\in\{A,B\}$ and $o\in\{s,f\}$,
\begin{equation}\label{eq:layer-increment}
n(T_j^o z)=n(z)+1.
\end{equation}
Starting from $z_0$, exactly one seller is selected and one outcome is
observed each period. Hence $n$ counts completed calendar periods: before
period $t$, $n=t-1$. In contrast, $n_j$ counts only periods in which
seller $j$ was selected. A large $n=n_A+n_B$ need not imply that both
individual counts are large.

When $B$ is selected, $A$'s own counts stay fixed, but the joint state
changes to $T_B^s z$ or $T_B^f z$. These changes can affect both future
demand and both sellers' subsequent product choices. Therefore $A$'s
unselected branch cannot be written as $\gamma D_B(z)V_A(z)$.
There is no unchanged-joint-state continuation term. The original
outside-option denominator $1-\gamma[1-D(S,F)]$ came from rearranging
such a self-loop and has no corresponding justification here.

\section{Policies and fixed-policy evaluation}

\subsection{Product policies and actual outcome probabilities}

A stationary Markov product policy for seller $j$ is a function
$\sigma_j:\mathbb N_0^4\to[0,1]$, with
\begin{equation}\label{eq:policy}
\sigma_j(z)=\Pb(i_t=2\mid y_t=j,z).
\end{equation}
This probability is the seller's endogenous product policy, conditional
on being selected. It is neither demand nor a customer belief.
A deterministic policy is the special case $\sigma_j(z)\in\{0,1\}$
at every state; it need not prescribe the same product at different states.

Under this policy, the actual success probability conditional on selection
and the expected product cost are respectively
\begin{equation}\label{eq:policy-outcomes}
\begin{aligned}
q_j(z)&=(1-\sigma_j(z))p_1+\sigma_j(z)p_2
       =p_1+\sigma_j(z)(p_2-p_1),\\
\bar c_j(z)&=(1-\sigma_j(z))c_1+\sigma_j(z)c_2
       =c_1+\sigma_j(z)(c_2-c_1).
\end{aligned}
\end{equation}
These are averages over the selected seller's product randomization.
Conditional on its realized product $i$, success still occurs with
probability $p_i$. In particular, $q_j(z)$ is not $\theta_j$, a Thompson
sample, or a posterior mean.

Let $\sigma=(\sigma_A,\sigma_B)$ be a fixed pair of policies. The
transition probabilities before customer selection are
\begin{equation}\label{eq:transition-probabilities}
\Pb(z'=T_k^s z\mid z,\sigma)=D_k(z)q_k(z),\qquad
\Pb(z'=T_k^f z\mid z,\sigma)=D_k(z)[1-q_k(z)],
\end{equation}
for $k\in\{A,B\}$, where $z'$ is next period's state. The four
probabilities sum to $D_A(z)+D_B(z)=1$.

\subsection{Discounted profits and the evaluation equation}

For each $j\in\{A,B\}$, the objective is to maximize
\begin{equation}\label{eq:objective}
\E\left[\sum_{t=1}^{\infty}\gamma^{t-1}
\Ind{y_t=j}(R-c_{i_t})\right]
\end{equation}
by choosing its own product strategy, taking the customer's rule and
the other seller's strategy as given. Discounting is by calendar period,
including periods when the other seller serves the customer.

For a fixed pair $\sigma$, define its evaluation value before customer
selection by
\begin{equation}\label{eq:fixed-value}
U_j^\sigma(z)=\E_z^\sigma\left[
\sum_{t=1}^{\infty}\gamma^{t-1}\Ind{y_t=j}(R-c_{i_t})\right].
\end{equation}
The subscript $z$ starts the continuation problem at state $z$; the
superscript specifies that both policies are followed thereafter.
The expectation averages over customer samples, product randomization,
and actual outcomes. Since period profit lies between zero and $R-c_1$,
\begin{equation}\label{eq:value-bound}
0\leq U_j^\sigma(z)\leq\frac{R-c_1}{1-\gamma}.
\end{equation}

Separating current profit from next period's value and conditioning on
the selected seller gives
\begin{equation}\label{eq:policy-evaluation}
\boxed{\begin{aligned}
U_j^\sigma(z)={}&D_j(z)[R-\bar c_j(z)]\\
&+\gamma\sum_{k\in\{A,B\}}D_k(z)
\left[q_k(z)U_j^\sigma(T_k^s z)
+(1-q_k(z))U_j^\sigma(T_k^f z)\right].
\end{aligned}}
\end{equation}
If $j$ is selected, it earns its expected current margin and continues
after an update to its own counts. If the other seller is selected,
$j$ earns zero now, but its continuation value changes with that
seller's outcome. Both branches are discounted by $\gamma$ because
both consume one calendar period. In every continuation term the value
belongs to seller $j$, even when the outcome comes from seller $k\ne j$.

\section{Best responses and equilibrium product choice}

\subsection{Best response to a fixed opposing policy}

Hold $B$'s entire function $\sigma_B$ fixed. Write $V_A^{\sigma_B}(z)$
for $A$'s optimal value before selection when it can choose its product
strategy at every future opportunity. In contrast to $U_A^\sigma$,
this value optimizes $A$'s policy. Formally,
\begin{equation}\label{eq:best-response-value}
V_j^{\sigma_{-j}}(z)=\sup_{\tau_j}
\E_z^{\tau_j,\sigma_{-j}}\left[
\sum_{t=1}^{\infty}\gamma^{t-1}\Ind{y_t=j}(R-c_{i_t})\right],
\end{equation}
where $-j$ denotes the other seller and $\tau_j$ ranges over admissible
product strategies for $j$, including history-dependent deviations.
With the opposing Markov policy fixed, the public state contains all
information relevant to rewards and future transitions, so the
best-response problem admits the following Markov Bellman formulation.

If $A$ is selected, it compares its two current products using its own
optimal continuation value. If $B$ is selected, $B$ uses the fixed
$\sigma_B(z)$, with success probability $q_B(z)$. Consequently,
\begin{equation}\label{eq:br-a}
\begin{aligned}
V_A^{\sigma_B}(z)={}&D_A(z)\max_{i\in\{1,2\}}
\Bigl\{R-c_i+\gamma\bigl[p_iV_A^{\sigma_B}(T_A^s z)
+(1-p_i)V_A^{\sigma_B}(T_A^f z)\bigr]\Bigr\}\\
&+\gamma D_B(z)\bigl[q_B(z)V_A^{\sigma_B}(T_B^s z)
+(1-q_B(z))V_A^{\sigma_B}(T_B^f z)\bigr].
\end{aligned}
\end{equation}
Analogously, holding $\sigma_A$ fixed gives
\begin{equation}\label{eq:br-b}
\begin{aligned}
V_B^{\sigma_A}(z)={}&D_B(z)\max_{i\in\{1,2\}}
\Bigl\{R-c_i+\gamma\bigl[p_iV_B^{\sigma_A}(T_B^s z)
+(1-p_i)V_B^{\sigma_A}(T_B^f z)\bigr]\Bigr\}\\
&+\gamma D_A(z)\bigl[q_A(z)V_B^{\sigma_A}(T_A^s z)
+(1-q_A(z))V_B^{\sigma_A}(T_A^f z)\bigr].
\end{aligned}
\end{equation}
The maximum is inside the selected seller's branch because product
choice occurs only after selection. Randomizing cannot improve on the
larger of the two conditional values; it is optimal only at a tie.

When $A$ deviates, the function $\sigma_B$ remains fixed, but $B$ applies
it to whichever states the deviation generates. For example,
$\sigma_B(T_A^s z)$ and $\sigma_B(T_A^f z)$ can differ. Neither $B$'s
realized future products nor the future state trajectory is frozen.

For a fixed opposing policy, let $\mathcal T_j^{\sigma_{-j}}$ denote
the Bellman operator on the right-hand side of the corresponding
equation. For any bounded functions $v,w$, let
$\|v-w\|_\infty=\sup_z|v(z)-w(z)|$. Probability weights sum to one,
and taking a maximum cannot increase the largest difference between
the two product values, so
\begin{equation}\label{eq:contraction}
\|\mathcal T_j^{\sigma_{-j}}v-\mathcal T_j^{\sigma_{-j}}w\|_\infty
\leq\gamma\|v-w\|_\infty.
\end{equation}
The bounded-function space is complete, hence this contraction gives
a unique bounded best-response value for each fixed opposing policy.
The finite product set permits a maximizing stationary selector.
The same probability argument gives a unique bounded evaluation value
for each fixed pair in \eqref{eq:policy-evaluation}. These statements
do not imply uniqueness of optimal products at ties or uniqueness of
the coupled game equilibrium.

\subsection{The future value difference and product choice}

For a candidate equilibrium pair $\sigma$, abbreviate its evaluated
values as $V_j=U_j^\sigma$; optimality of these candidate values will
be imposed below. As in the paper, define
\begin{equation}\label{eq:increments}
\Delta p=p_2-p_1>0,\qquad \Delta c=c_2-c_1>0.
\end{equation}
Conditional on seller $j$ being selected, choosing product $i$ now and
following the candidate policies thereafter has value
\begin{equation}\label{eq:conditional-q}
Q_j^i(z)=R-c_i+\gamma\left[
p_iV_j(T_j^s z)+(1-p_i)V_j(T_j^f z)\right].
\end{equation}
The conditional gain $a_j(z)$ from product $2$ rather than product $1$
is obtained by explicit subtraction:
\begin{equation}\label{eq:gain}
\begin{aligned}
a_j(z)&=Q_j^2(z)-Q_j^1(z)\\
&=\left\{R-c_2+\gamma[p_2V_j(T_j^s z)+(1-p_2)V_j(T_j^f z)]\right\}\\
&\quad-\left\{R-c_1+\gamma[p_1V_j(T_j^s z)+(1-p_1)V_j(T_j^f z)]\right\}\\
&=-(c_2-c_1)+\gamma\left[(p_2-p_1)V_j(T_j^s z)
+(p_1-p_2)V_j(T_j^f z)\right]\\
&=-\Delta c+\gamma\Delta p
\left[V_j(T_j^s z)-V_j(T_j^f z)\right].
\end{aligned}
\end{equation}
Product $2$ sacrifices $\Delta c$ in current margin and shifts
probability $\Delta p$ from the failure successor to the success
successor. The bracket is the future value difference for seller $j$.
Retaining the paper's discounting convention, its continuation gap is
\begin{equation}\label{eq:gap}
G_j(z)=\gamma\left[V_j(T_j^s z)-V_j(T_j^f z)\right],
\qquad a_j(z)=-\Delta c+\Delta p\,G_j(z).
\end{equation}
Both continuation values allow the two sellers to respond to the new
joint state. The formula by itself establishes neither a sign nor a
monotonicity property for this future value difference.

To see the cancellation before selection explicitly, define $C_j(z)$
as seller $j$'s discounted continuation contribution when the other
seller is selected:
\begin{equation*}
C_j(z)=\gamma D_{-j}(z)\left[
q_{-j}(z)V_j(T_{-j}^s z)+(1-q_{-j}(z))V_j(T_{-j}^f z)\right].
\end{equation*}
The two before-selection values for a one-period product deviation
differ by
\begin{equation}\label{eq:unselected-cancellation}
\bigl[D_j(z)Q_j^2(z)+C_j(z)\bigr]
-\bigl[D_j(z)Q_j^1(z)+C_j(z)\bigr]=D_j(z)a_j(z).
\end{equation}
The unselected branch cancels because seller $j$ makes no product
choice in that branch. Current demand is identical under the two
choices because selection precedes the hidden product choice. Since
$D_j(z)>0$, the conditional and before-selection gains have the same
sign. Demand still matters for the levels of all continuation values.

\subsection{Stationary Markov perfect equilibrium}

A stationary Markov perfect equilibrium is a policy pair
$\sigma=(\sigma_A,\sigma_B)$ that is mutually optimal at every state:
each seller's policy is a best response to the other's entire policy
function, including at states off the equilibrium path. Equivalently,
\begin{equation}\label{eq:mpe-values}
U_A^\sigma(z)=V_A^{\sigma_B}(z),\qquad
U_B^\sigma(z)=V_B^{\sigma_A}(z)\qquad\text{for every }z.
\end{equation}
``Stationary'' means the policy depends on the current state rather
than separately on calendar time. It does not mean constant products,
constant beliefs, or constant market shares. In this cumulative-count
model the state also records elapsed time through $n(z)$.

Using the evaluated candidate values in \eqref{eq:gain}, the equilibrium
product-choice conditions for each $j\in\{A,B\}$ are
\begin{equation}\label{eq:equilibrium-policy}
\sigma_j(z)=
\begin{cases}
1,&a_j(z)>0,\\
0,&a_j(z)<0,\\
\text{any number in }[0,1],&a_j(z)=0.
\end{cases}
\end{equation}
In the paper's terminology, product $2$ is optimal exactly when
$G_j(z)\geq\Delta c/\Delta p$, with strict preference when the
inequality is strict. This is a product-choice condition, not an
assumption that the equilibrium policy has a threshold in a count or
a posterior mean.

Conditions \eqref{eq:equilibrium-policy} must hold together with
\eqref{eq:policy-evaluation} for \emph{both} sellers at \emph{every}
state. They are sufficient as well as necessary: if evaluation and
the maximizing product conditions hold, each evaluated value satisfies
its own best-response Bellman equation. Uniqueness of that fixed-
opponent value then gives \eqref{eq:mpe-values}. At indifference,
any mixture is locally maximizing for the given continuation values;
changing a mixture in a policy pair still requires consistent
reevaluation of both sellers' values.

The sellers have symmetric primitives and initial beliefs, but their
realized counts need not be equal. No equality of reputations, constant
policy, symmetry restriction on equilibrium, or equilibrium uniqueness
is imposed. The separate best-response contractions do not establish
uniqueness of a mutually optimal policy pair.

\section{An exact finite-horizon benchmark}

For a separate finite-horizon game, let the positive integer $H$ be the
total number of periods from $z_0$. Write $V_j^H(z)$ for seller $j$'s
equilibrium value before selection with $H-n(z)$ periods remaining,
discounted relative to the current period. The horizon $H$ is fixed
throughout the recursion. Set
\begin{equation}\label{eq:terminal}
V_j^H(z)=0\quad\text{when }n(z)=H,\qquad j\in\{A,B\}.
\end{equation}
Work backwards through layers $n=H-1,H-2,\ldots,0$. At any state on
layer $n$, all four successors lie on layer $n+1$, so both sellers'
continuation values there are already determined.

Conditional on seller $j$ being selected, define
\begin{equation}\label{eq:finite-q}
Q_j^{H,i}(z)=R-c_i+\gamma\left[
p_iV_j^H(T_j^s z)+(1-p_i)V_j^H(T_j^f z)\right],
\qquad i\in\{1,2\}.
\end{equation}
For $A$ and $B$ separately, choose the product rule
\begin{equation}\label{eq:finite-policy}
i_j^H(z)\in\argmax_{i\in\{1,2\}}Q_j^{H,i}(z),\qquad
i_j^H(z)=1\ \text{if }Q_j^{H,1}(z)=Q_j^{H,2}(z).
\end{equation}
Having chosen \emph{both} products, calculate \emph{both} current
values using those same choices:
\begin{equation}\label{eq:finite-evaluation}
\begin{aligned}
V_j^H(z)={}&D_j(z)[R-c_{i_j^H(z)}]\\
&+\gamma\sum_{k\in\{A,B\}}D_k(z)
\left[p_{i_k^H(z)}V_j^H(T_k^s z)
+(1-p_{i_k^H(z)})V_j^H(T_k^f z)\right],
\quad j\in\{A,B\}.
\end{aligned}
\end{equation}
Thus $B$'s selected product affects $A$'s current value through the
branch where $B$ serves the customer, and conversely for $A$'s product
and $B$'s value. Each seller maximizes its own conditional value,
not the other seller's or their sum.

This recursion constructs a pure-strategy Markov equilibrium by backward
induction. At the terminal layer there are no further decisions.
Given equilibrium continuation play on the next layer, only the
selected seller acts at the current state, and its maximizing product
is optimal conditional on selection. The other seller has no current
decision in that branch. Applying this argument at every state and
layer gives sequentially optimal play throughout the finite game.
With continuation values fixed, a generic simultaneous $2$-by-$2$
Nash solver is therefore unnecessary. Choosing product $1$ at exact
indifference makes the recursion definite and selects an equilibrium;
it does not establish uniqueness of equilibrium in the game.

In the final decision period, $n(z)=H-1$, both successor values in
\eqref{eq:finite-q} vanish. Hence, for each seller,
\begin{equation}\label{eq:final-product}
Q_j^{H,2}(z)-Q_j^{H,1}(z)
=(R-c_2)-(R-c_1)=-\Delta c<0.
\end{equation}
Both sellers therefore prescribe product $1$ whenever selected, and
$V_j^H(z)=D_j(z)(R-c_1)$ on this layer. This is a terminal-period
result for the finite-horizon game. It does not establish eventual
extinction of product $2$ in the original infinite-horizon model.
No convergence of finite-horizon equilibrium policies to a particular
infinite-horizon equilibrium is asserted here.

\section{What is established and what remains open}

The belief rule and timing directly imply the demand integrals, the
four state transitions, the fixed-policy and best-response equations,
and the product-choice condition in terms of the future value difference.
Backward induction gives the stated finite-horizon equilibrium and its
final-period use of product $1$. These are analytical consequences
of the specified model.

Characterizing the infinite-horizon equilibrium policy pair, the shape
of its product-choice regions, possible multiplicity, and long-run
product use or market shares requires further equilibrium analysis.
In particular, the present derivation establishes neither threshold
policies nor extinction or persistence of product $2$. Numerical
analysis, when undertaken separately, may investigate these questions;
it would not by itself turn observed patterns into general results.

\end{document}
